Hubo que adaptar la receta a una máquina más pequeña. Cuatro cambios, cada uno con su razón.

Los hilos: dividimos la lista del paper entre seis. Así se conserva la proporción: 24 hilos es toda la máquina, como 144 allí.

La duración: cinco rondas de dos segundos en vez de siete de diez. Cada punto en un proceso nuevo y en orden aleatorio. Pasamos de cinco horas a una.

La latencia: medimos con 1, 24 y 56 hilos. Y medimos también a Williams, que el paper deja fuera. Al hacerlo encon- tramos un problema de compilación que veremos en la parte 7.

Recortar tiene un precio, y lo medimos: está en la parte 10.
